\documentclass[11pt,a4paper]{article}
\usepackage[T1]{fontenc}
\usepackage[utf8]{inputenc}
\usepackage{lmodern}
\usepackage[margin=25mm]{geometry}
\usepackage{amsmath,amsthm,amssymb,mathtools}
\usepackage{microtype}
\usepackage{booktabs,longtable,array,tabularx}
\usepackage{enumitem}
\usepackage{xcolor}
\usepackage{listings}
\usepackage[colorlinks=true,linkcolor=blue!40!black,citecolor=blue!40!black,urlcolor=blue!50!black]{hyperref}
\hypersetup{pdftitle={Expressing algebraic numbers by radicals: a complete replacement for ToRadicals}}

\lstset{basicstyle=\ttfamily\small,columns=fullflexible,keepspaces=true,breaklines=true,
  frame=single,framerule=0.3pt,rulecolor=\color{black!30},xleftmargin=4pt,xrightmargin=4pt,
  aboveskip=6pt,belowskip=6pt,showstringspaces=false}

\newtheorem{theorem}{Theorem}[section]
\newtheorem{proposition}[theorem]{Proposition}
\newtheorem{lemma}[theorem]{Lemma}
\newtheorem{corollary}[theorem]{Corollary}
\theoremstyle{definition}
\newtheorem{definition}[theorem]{Definition}
\newtheorem{example}[theorem]{Example}
\newtheorem{algorithm}[theorem]{Algorithm}
\theoremstyle{remark}
\newtheorem{remark}[theorem]{Remark}

\newcommand{\Q}{\mathbb{Q}}
\newcommand{\Z}{\mathbb{Z}}
\newcommand{\C}{\mathbb{C}}
\newcommand{\R}{\mathbb{R}}
\newcommand{\F}{\mathbb{F}}
\newcommand{\Qbar}{\overline{\Q}}
\newcommand{\Gal}{\operatorname{Gal}}
\newcommand{\Tr}{\operatorname{Tr}}
\newcommand{\Res}{\operatorname{Res}}
\newcommand{\Stab}{\operatorname{Stab}}
\newcommand{\AGL}{\operatorname{AGL}}
\newcommand{\rad}{\operatorname{rad}}
\newcommand{\Frob}{\operatorname{Frob}}
\newcommand{\code}[1]{\texttt{#1}}
\newcommand{\Root}{\code{Root}}
\newcommand{\ToRadicals}{\code{ToRadicals}}

\title{Expressing algebraic numbers by radicals:\\ a complete replacement for \ToRadicals}
\author{Prepared for Vladimir Reshetnikov}
\date{September 2026}

\begin{document}
\maketitle

\begin{abstract}
\noindent
Mathematica StackExchange question 34011 asks why \ToRadicals{} cannot rewrite every
\Root{} object that has a radical expression, whether the obstruction is undecidability
or cost, and whether a complete replacement can be written.  We answer all three
questions.  A radical expression exists exactly when the Galois group of the minimal
polynomial is solvable; this is decidable, and even in polynomial time, so the built-in
limitation is a matter of scope and cost, not of principle.  We give a complete
algorithm: cheap structural recognizers that produce short formulas (functional
decomposition, generalized reciprocal symmetry, Dickson polynomials, and the pair-sum
resolvent that systematizes the search in the notebook attached to the question),
and a general Galois--Kummer descent that is complete for every solvable input and
proves nonsolvability otherwise.  Every step is proved.  Both examples of the question
are worked out by hand and by machine, with a proof that the branch returned is the
requested one.  Two implementations are provided and tested: a Wolfram Language package
and a Python port on python-flint, both reusing the splitting-field engine of the
companion root-decomposition project.  Every claim about running code in this article
was executed; the measurements are reported as observed.
\end{abstract}

\tableofcontents

\section{Introduction}\label{sec:intro}

The documentation of \ToRadicals{} states that there are cases in which an expression in
radicals exists but the function cannot find it.  The question gives two such cases,
\begin{equation}\label{eq:examples}
\alpha = \Root[x^6+x^4-x^3-x^2-1,\,2], \qquad
\beta  = \Root[5x^5-25x^3+25x+6,\,5],
\end{equation}
with the closed forms
\begin{equation}\label{eq:closedforms}
\begin{aligned}
\alpha&=\frac{1}{\sqrt[3]{36}}\Bigl(\frac{3}{\sqrt{\gamma}}+\frac{\gamma}{2}\Bigr),\qquad
\gamma=\sqrt[3]{2\delta}-8\sqrt[3]{3/\delta},\qquad \delta=9+\sqrt{849},\\
\beta&=\sqrt[5]{\tfrac{-3+4i}{5}}+\sqrt[5]{\tfrac{-3-4i}{5}},
\end{aligned}
\end{equation}
and asks: what is the nature of the restriction, is the problem undecidable or merely
expensive, and can one write a version of \ToRadicals{} that succeeds whenever a radical
expression exists?

\subsection{Summary of results}

\begin{enumerate}[nosep]\sloppy
\item \textbf{Nature of the restriction.}  An algebraic number has a radical expression
if and only if the Galois group of its minimal polynomial is solvable
(Theorem~\ref{thm:galois}).  Solvability of a Galois group is decidable, and Landau and
Miller showed in 1985 that it is decidable in polynomial time~\cite{LandauMiller}.  The
built-in function is documented to be complete through degree four and to recognize
some special forms; on the two examples above and on the cyclic quintic
$x^5+x^4-4x^3-3x^2+3x+1$ (whose roots are $2\cos(2\pi k/11)$) it returns the \Root{}
object unchanged (Section~\ref{sec:builtin}).  The obstruction is therefore neither
undecidability nor the absence of an algorithm; it is that the general construction
requires the splitting field, whose degree can be as large as the group order.
\item \textbf{A complete algorithm.}  Section~\ref{sec:descent} gives the general
Galois--Kummer descent with proofs and a completeness theorem
(Theorem~\ref{thm:complete}).  Section~\ref{sec:structural} gives four structural
reductions that produce short formulas without the splitting field, each with a proof;
the two examples of the question are instances of two of them
(Section~\ref{sec:examples}).  Section~\ref{sec:negative} gives cheap proofs of
nonsolvability from Frobenius cycle types.
\item \textbf{Two implementations.}  \code{RootToRadicals.wl} (Wolfram Language) and
\code{roottoradicals.py} (Python, python-flint and SymPy) implement the same
algorithms; both were executed on the examples and on regression suites
(Section~\ref{sec:impl}).  Both examples of the question are solved in well under a
second by the structural layer, and in 41 and 18 seconds (Wolfram) or 24 and 12 seconds
(Python) when the general descent is forced.
\item \textbf{The notebook and the nine reports.}  Section~\ref{sec:notebook} explains
what the notebook attached to the question computes and which theorem replaces each of
its devices.  Section~\ref{sec:reports} compares the nine independently written reports
that preceded this article.
\end{enumerate}

\subsection{Conventions}\label{sec:conventions}

A \emph{radical expression} is built from rational numbers, the imaginary unit $i$,
addition, multiplication, and powers with rational exponents, where $z^{p/q}$ denotes
the principal branch $\exp\bigl(\tfrac{p}{q}\log z\bigr)$ with $-\pi<\operatorname{Im}\log z\le\pi$.
This is exactly the grammar of Wolfram's \code{Plus}, \code{Times}, \code{Power} with
rational exponents and \code{I}; in particular $\zeta_m=(-1)^{2/m}=e^{2\pi i/m}$ is a
radical expression, so roots of unity need no special treatment.  The \emph{depth} of a
radical expression is the nesting depth of its non-integer powers.  Throughout,
$p\in\Z[x]$ is the minimal polynomial of the input, $n=\deg p$, $L$ is its splitting
field in $\C$, and $G=\Gal(L/\Q)$, viewed as a permutation group on the $n$ roots.
\Root{}$[p,k]$ denotes the $k$-th root in Wolfram's ordering (real roots first,
increasing; then non-real roots in conjugate pairs).

\section{The problem}\label{sec:problem}

There are three distinct tasks.
\begin{enumerate}[nosep]
\item \emph{Decide} whether some radical expression equals the given number.
\item \emph{Construct} one.
\item \emph{Identify the requested conjugate}: the formula must equal the root with the
given index, not merely some root of $p$.
\end{enumerate}
The third task is not a formality.  A formula in radicals with principal branches
determines a single complex number, and a change of one branch usually gives a different
conjugate; the notebook of the question spends most of its effort on exactly this point
(Section~\ref{sec:notebook}).  Our algorithms select branches by a certified numerical
comparison among finitely many separated candidates and then verify the final identity
exactly.

The criterion must be applied to the \emph{minimal} polynomial of the selected number.
If the user supplies a reducible or non-minimal annihilating polynomial, its Galois
group may be larger (even nonsolvable) although the selected root lies in a solvable
factor.  All functions below therefore start by computing the minimal polynomial.

\section{Why the built-in function fails, and why the problem is decidable}\label{sec:builtin}

Wolfram's \ToRadicals{} is documented to convert \Root{} objects of degree at most four,
and \Root{} objects of some recognized special forms (for instance binomials
$x^m-c$ and polynomials that \code{Decompose} reduces to such pieces).  In Wolfram
15.0.1 we measured:

\begin{center}
\begin{tabular}{@{}lcl@{}}
\toprule
input & \ToRadicals & remark\\
\midrule
$\Root[x^6+x^4-x^3-x^2-1,2]$ & unchanged & example 1 of the question\\
$\Root[5x^5-25x^3+25x+6,5]$ & unchanged & example 2 of the question\\
$\Root[x^5+x^4-4x^3-3x^2+3x+1,1]$ & unchanged & cyclic quintic, $2\cos(2\pi/11)$\\
$\Root[x^3-3x+1,1]$ & radicals & degree $\le 4$\\
$\Root[x^4-x-1,1]$ & radicals & degree $\le 4$\\
$\Root[x^5-2,1]$, $\Root[x^8+1,1]$ & radicals & binomials\\
$\Root[x^6-2x^3-1,1]$ & radicals & decomposable\\
\bottomrule
\end{tabular}
\end{center}

None of the three failures is caused by undecidability.  The Galois group of a
polynomial is a finite computable object (by any of the classical resolvent methods, or
by the numerical resolvent method of our engine, Section~\ref{sec:engine}), and
solvability of a finite group is decided by iterating the commutator subgroup.
Landau and Miller~\cite{LandauMiller} proved that solvability by radicals can be decided
in time polynomial in the degree and the coefficient size, without computing the Galois
group; the \emph{construction} of an expression may however require the splitting field,
whose degree $[L:\Q]=|G|$ can be $n!$.  This is the real cost, and it is why a
general-purpose library function stops at degree four: beyond it there is no bound on
the size of the answer that is polynomial in the input.

\section{The criterion}\label{sec:criterion}

\begin{theorem}[Galois]\label{thm:galois}
Let $a\in\Qbar$ with minimal polynomial $p$, splitting field $L$ and $G=\Gal(L/\Q)$.
Then $a$ is equal to a radical expression if and only if $G$ is solvable.
\end{theorem}

\begin{proof}
See \cite[Chapter 5]{Milne} or \cite[\S 14.7]{DummitFoote}.  For the direction we use
constructively, Section~\ref{sec:descent} gives an explicit procedure.
\end{proof}

The constructive direction needs roots of unity in the base field.  The following
reduction is used by every algorithm in this article.

\begin{lemma}[Cyclotomic base]\label{lem:base}
Let $m=\rad|G|$ be the product of the distinct primes dividing $|G|$, $B=\Q(\zeta_m)$,
and $M=LB$ the compositum.  Then
\begin{enumerate}[nosep]
\item[(i)] $H:=\Gal(M/B)\cong\Gal(L/L\cap B)$, a subgroup of $G$;
\item[(ii)] $H$ is solvable if and only if $G$ is solvable;
\item[(iii)] every prime dividing $|H|$ divides $m$, so $B$ contains $\zeta_q$ for every
composition factor $q$ of $H$.
\end{enumerate}
\end{lemma}

\begin{proof}
(i) is the theorem on natural irrationalities \cite[Chapter 3]{Milne}: restriction
$\Gal(M/B)\to\Gal(L/L\cap B)$ is an isomorphism because $M=LB$ and $L/\Q$ is Galois.
(ii) A subgroup of a solvable group is solvable; conversely $G/\Gal(L/L\cap B)\cong\Gal(L\cap B/\Q)$
is a quotient of the abelian group $\Gal(B/\Q)$, and an extension of a solvable group by
an abelian group is solvable.  (iii) follows from (i) and Lagrange's theorem.
\end{proof}

The lemma converts the problem into a tower of cyclic extensions of prime degree $q$ over
fields containing $\zeta_q$, which is the setting of Kummer theory
(Section~\ref{sec:descent}).  Note that $m$ omits the prime $2$ in practice: $\zeta_2=-1$
is rational.

\section{Structural reductions}\label{sec:structural}

The reductions of this section find short formulas from the shape of $p$ alone, without
computing the splitting field.  Each is a theorem of the form ``if $p$ has this shape,
then $a$ is obtained from a number of smaller degree by one explicit radical step''.
The implementations try them in the order given, recursing on the smaller number, and
fall back to the general descent.

\subsection{Functional decomposition}

\begin{proposition}\label{prop:decompose}
Let $p=g\circ h$ with $g,h\in\Q[x]$ of degrees $\ge2$, and let $a$ be a root of $p$.
Then $v=h(a)$ is a root of $g$, and $a$ is a root of $h(x)-v\in\Q(v)[x]$.  If $v$ has a
radical expression, then so does $a$ whenever $\deg h\le4$ or $h(x)=cx^k+d$ is a
binomial.
\end{proposition}

\begin{proof}
$g(v)=g(h(a))=p(a)=0$.  For $\deg h\le4$ the roots of $h(x)-v$ are given by the
classical formulas with coefficients in $\Q(v)$; for a binomial they are
$\zeta_k^{e}\bigl((v-d)/c\bigr)^{1/k}$.  Both are radical expressions in a radical
expression of $v$.
\end{proof}

\code{Decompose} returns a complete decomposition $p=g_1\circ g_2\circ\dots\circ g_r$
(Ritt's theorem guarantees that the multiset of degrees is well defined).  The
implementation peels from the outside: $v_1$ is a root of $g_1$, then $v_2$ is a root of
$g_2(x)-v_1$, and so on; at each stage the exact value $v_i=(g_{i+1}\circ\dots\circ g_r)(a)$
is known (it is an algebraic number computed by \code{RootReduce} or by a resultant), and
the branch of the radical formula is the candidate that matches this exact value
numerically.  The chain fails only when some inner piece of degree $>4$ is not a
binomial; the general descent then takes over.

\subsection{Generalized reciprocal symmetry}\label{sec:reciprocal}

\begin{proposition}\label{prop:reciprocal}
Let $n=2m$ and suppose $p(x)=x^mP(x+c/x)$ with $c\in\Q^\times$ and $P\in\Q[x]$ of
degree $m$.  If $a$ is a root of $p$, then $y=a+c/a$ is a root of $P$, and
$a=\tfrac12\bigl(y\pm\sqrt{y^2-4c}\bigr)$; the other sign gives the conjugate $c/a$.
\end{proposition}

\begin{proof}
The constant term of the monic form of $p$ is $c^m\neq0$, so $a\neq0$ and
$0=p(a)=a^mP(a+c/a)$ gives $P(y)=0$.  The numbers $a$ and $c/a$ are the two roots of
$x^2-yx+c$.
\end{proof}

\emph{Detection.}  Write $p$ monic.  Then $c^m$ is the constant term, so $c$ is a
rational $m$-th root of it (two candidates $\pm c$ when $m$ is even); for each candidate,
subtract $\operatorname{coef}_{m+k}(p)\,x^m(x+c/x)^k$ for $k=m,m-1,\dots,0$ and check that
the remainder vanishes.  This is a finite exact test.

\emph{Branch coupling.}  The formula carries $c/a$ as the other root of the same
quadratic.  Extracting $a$ and $c/a$ by two independent radical formulas and then
choosing branches for each would not be sound: the two choices are coupled by
$a\cdot(c/a)=c$.  The same remark applies to Dickson polynomials below, where the
term $c/u$ must be formed from the \emph{same} $u$.  The first example of the question
(Section~\ref{sec:examples}) is of this type with $c=-1$.

\subsection{Dickson polynomials}\label{sec:dickson}

The Dickson polynomial $D_k(x,c)$ is defined by $D_0=2$, $D_1=x$,
$D_k=xD_{k-1}-cD_{k-2}$, and satisfies the functional equation
\begin{equation}\label{eq:dickson}
D_k\Bigl(u+\frac{c}{u},\,c\Bigr)=u^k+\frac{c^k}{u^k}.
\end{equation}

\begin{proposition}\label{prop:dickson}
Suppose that after the translation $x\mapsto x+t$ that kills the $x^{n-1}$ term, the monic
form of $p$ becomes $D_n(x,c)-b$ with $c\in\Q^\times$, $b\in\Q$.  Let $s$ be a root of
$s^2-bs+c^n=0$ and $u=s^{1/n}$.  Then the $n$ numbers
\[
t+\zeta_n^{\,j}u+\frac{c}{\zeta_n^{\,j}u},\qquad j=0,\dots,n-1,
\]
are roots of $p$; if they are distinct they are all of them.
\end{proposition}

\begin{proof}
With $s'=c^n/s$ we have $s+s'=b$.  By \eqref{eq:dickson},
$D_n(\zeta^ju+c/(\zeta^ju),c)=(\zeta^ju)^n+c^n/(\zeta^ju)^n=s+s'=b$.
\end{proof}

\emph{Detection.}  $D_n(x,c)=x^n-ncx^{n-2}+\dots$, so $c$ is read off the $x^{n-2}$
coefficient of the translated polynomial and the identity $p=D_n(x,c)-b$ is checked
exactly.  The second example of the question is $5\bigl(D_5(x,1)+\tfrac65\bigr)$.

\subsection{The pair-sum resolvent}\label{sec:pairsum}

The following identity is classical.

\begin{lemma}\label{lem:pairsum}
Let $p$ have leading coefficient $\ell$ and roots $a_1,\dots,a_n$, and let
$R_2(y)=\prod_{i<j}(y-a_i-a_j)$.  Then
\[
\Res_x\bigl(p(x),\,p(y-x)\bigr)=2^n\ell^{2n-1}\,p(y/2)\,R_2(y)^2 .
\]
\end{lemma}

\begin{proof}
$\Res_x(p(x),p(y-x))=\ell^n\prod_i p(y-a_i)=\ell^{2n}\prod_{i,j}(y-a_i-a_j)$.  The
diagonal terms give $\prod_i(y-2a_i)=2^n\ell^{-1}p(y/2)$; the off-diagonal terms give
$R_2(y)^2$.
\end{proof}

So $R_2$ is computed exactly by a resultant, an exact division, and a square root of a
factorization.  Its roots are the pair sums $a_i+a_j$.

\begin{proposition}\label{prop:pairsum}
Let $y_0=a+a'$ be a pair sum with $a'$ a conjugate of $a$, and let $K=\Q(y_0)$.  Factor
$p$ over $K$ and let $f\in K[x]$ be the irreducible factor with $f(a)=0$.  If $y_0$ has a
radical expression and $\deg f\le4$ (or $f$ is a binomial over $K$), then $a$ has a
radical expression, obtained by substituting the expression of $y_0$ into the
coefficients of $f$ and solving.
\end{proposition}

\begin{proof}
Immediate from the classical formulas, as in Proposition~\ref{prop:decompose}.
\end{proof}

In practice one takes the factors of $R_2$ over $\Q$ of degree $2\le m<n$, in increasing
degree, tests numerically which of their roots are of the form $a+a'$, factors $p$ over
$\Q(y_0)$ (Wolfram's \code{Factor} with \code{Extension} is fast for the small degrees
involved) and recurses on $y_0$.  When $a'=c/a$ this reproduces
Proposition~\ref{prop:reciprocal}, but the pair-sum reduction is more general: the
product $aa'$ need not be rational, it is a polynomial in $y_0$ read off the factor
$f$.  Section~\ref{sec:notebook} shows that this is exactly what the notebook's
\code{findExtension} searches for by elimination.

\section{The two examples of the question}\label{sec:examples}

\subsection{The sextic}

$p(x)=x^6+x^4-x^3-x^2-1$ is monic with constant term $-1$, so in
Proposition~\ref{prop:reciprocal} $m=3$ and $c^3=-1$, $c=-1$.  Peeling gives
\[
p(x)=x^3\Bigl(\bigl(x-\tfrac1x\bigr)^3+4\bigl(x-\tfrac1x\bigr)-1\Bigr),
\qquad P(y)=y^3+4y-1 .
\]
Cardano's formula for the real root of $P$ (discriminant $-4\cdot4^3-27<0$, one real
root) gives
\[
y=u-\frac{4}{3u},\qquad u=\Bigl(\frac{9+\sqrt{849}}{18}\Bigr)^{1/3},
\]
and $\alpha=\tfrac12\bigl(y+\sqrt{y^2+4}\bigr)$.  This is the expression the package
returns (in Wolfram's own arrangement of the cube roots), and it is equal to the closed
form \eqref{eq:closedforms} of the question: with $\gamma$ as there,
$\gamma/\sqrt[3]{36}=y$ and $3/(\sqrt[3]{36}\sqrt{\gamma})=\tfrac12\sqrt{y^2+4}-\tfrac y2$,
as one checks by squaring; \code{RootReduce} confirms the identity.

\emph{Branch.}  $p$ has exactly two real roots, $\alpha_1<0<\alpha_2$, and
$\Root[p,2]=\alpha_2$ is the positive one.  The map $x\mapsto x-1/x$ is increasing on
$(0,\infty)$ and on $(-\infty,0)$, and $y$ is the unique real root of $P$; the two real
solutions of $x-1/x=y$ are $\tfrac12(y\pm\sqrt{y^2+4})$, of which the one with the plus
sign is positive.  Hence the formula equals $\Root[p,2]$; the minus sign gives
$\Root[p,1]=-1/\alpha_2$.  The four non-real roots come from the two non-real roots of
$P$, which are the other two branches of the cube root in $u$.

\subsection{The quintic}

$p(x)=5x^5-25x^3+25x+6=5\bigl(x^5-5x^3+5x+\tfrac65\bigr)$, and
$D_5(x,1)=x^5-5x^3+5x$, so $p=5\,(D_5(x,1)-b)$ with $b=-6/5$ and $c=1$.  The quadratic
$s^2+\tfrac65s+1=0$ has roots $s=\tfrac{-3\pm4i}{5}$, of modulus $1$.  By
Proposition~\ref{prop:dickson} the roots are $\zeta_5^ju+\zeta_5^{-j}u^{-1}$ with
$u=s^{1/5}$; since $|s|=1$, $u^{-1}=\bar u$ and the roots are $2\operatorname{Re}(\zeta_5^ju)=2\cos\bigl(\tfrac{\theta+2\pi j}{5}\bigr)$, $\theta=\arg s$, all
real.  The package returns $s^{1/5}+s^{-1/5}$ with $s=\tfrac{-3+4i}{5}$, which is
\eqref{eq:closedforms} because $s^{-1/5}=\bar s^{\,1/5}$ for $|s|=1$ (both are the
principal fifth root of $\bar s$: the principal argument of $\bar s$ is $-\theta$ and
$s^{-1/5}=e^{-i\theta/5}$).

\emph{Branch.}  $\theta=\arg\tfrac{-3+4i}{5}\in(\pi/2,\pi)$, so $\theta/5\in(\pi/10,\pi/5)$
and $2\cos(\theta/5)$ is the largest of the five values $2\cos((\theta+2\pi j)/5)$; all
five roots are real, so $\Root[p,5]$ is the largest root, which is the returned
branch $j=0$.  The other conjugates are $j=1,\dots,4$, and the package returns each of
them for the corresponding index.

\section{The general Galois--Kummer descent}\label{sec:descent}

\subsection{Setting}

Let $B=\Q(\zeta_m)$ and $M=LB$ as in Lemma~\ref{lem:base}, and $H=\Gal(M/B)$.  Choose a
composition series
\[
H=H_0\rhd H_1\rhd\dots\rhd H_r=1,\qquad |H_{i-1}/H_i|=q_i\text{ prime},
\]
which exists exactly when $H$ is solvable (the implementation builds it from the
commutator subgroup: $H_i$ is a subgroup of prime index in $H_{i-1}$ containing
$[H_{i-1},H_{i-1}]$, hence normal).  Let $F_i=M^{H_i}$ be the fixed fields, so
$B=F_0\subset F_1\subset\dots\subset F_r=M$ with $F_i/F_{i-1}$ cyclic of prime degree
$q_i$, generated by the image of any $\sigma_i\in H_{i-1}\setminus H_i$.

\subsection{One Kummer step}

\begin{proposition}[Lagrange resolvents]\label{prop:lagrange}
Let $F\subset E$ be fields with $E/F$ cyclic of prime degree $q$, generated by $\sigma$,
and suppose $\zeta=\zeta_q\in F$.  For $v\in E$ put
\[
R_k=\sum_{j=0}^{q-1}\zeta^{-kj}\sigma^j(v),\qquad k=0,\dots,q-1 .
\]
Then
\begin{enumerate}[nosep]
\item[(i)] $\sigma(R_k)=\zeta^kR_k$; hence $R_k^q\in F$, and $R_0\in F$;
\item[(ii)] $\sum_{k=0}^{q-1}R_k=q\,v$;
\item[(iii)] if $R_{k_1}\neq0$ for some $k_1\neq0$, then for every $k$ the element
$c_k=R_k/R_{k_1}^{\,m}$ with $m\equiv k\,k_1^{-1}\pmod q$ lies in $F$;
\item[(iv)] if $v\notin F$ then some $R_k$ with $k\neq0$ is nonzero.
\end{enumerate}
\end{proposition}

\begin{proof}
(i) $\sigma(R_k)=\sum_j\zeta^{-kj}\sigma^{j+1}(v)=\zeta^{k}\sum_{j'}\zeta^{-kj'}\sigma^{j'}(v)$,
using $\sigma^q=1$ and $\zeta^{q}=1$.  Then $\sigma(R_k^q)=\zeta^{kq}R_k^q=R_k^q$, so
$R_k^q$ is fixed by $\Gal(E/F)$.  (ii) $\sum_kR_k=\sum_j\bigl(\sum_k\zeta^{-kj}\bigr)\sigma^j(v)$
and the inner sum is $q$ for $j=0$ and $0$ otherwise.  (iii) $\sigma(c_k)=\zeta^kR_k/(\zeta^{k_1m}R_{k_1}^m)=c_k$
since $k_1m\equiv k$.  (iv) If all $R_k$ with $k\ne0$ vanish, (ii) gives $qv=R_0\in F$.
\end{proof}

Applied to $F=F_{i-1}$, $E=F_i$, $\sigma=\sigma_i$ and an element $v\in F_i$, the
proposition expresses $v$ through elements of $F_{i-1}$ and $q$-th roots of elements of
$F_{i-1}$ in one of two ways:
\begin{align}
\text{(Fourier form)}\qquad v&=\frac1q\Bigl(R_0+\sum_{k\ne0}\zeta^{e_k}\bigl(R_k^q\bigr)^{1/q}\Bigr),\label{eq:fourier}\\
\text{(eigenvector form)}\qquad v&=\frac1q\Bigl(R_0+u+\sum_{k\ne0,k_1}c_k\,u^{m_k}\Bigr),
\quad u=\zeta^{e}\bigl(R_{k_1}^q\bigr)^{1/q}.\label{eq:eigen}
\end{align}
The exponents $e_k,e$ are the \emph{branches}: $\bigl(R_k^q\bigr)^{1/q}$ is the principal
root, and exactly one of the $q$ numbers $\zeta^e(R_k^q)^{1/q}$ equals $R_k$.  The Fourier
form extracts $q-1$ radicals at this level; the eigenvector form extracts one and pushes
the ratios $c_k$, which are elements of $F_{i-1}$, one level down.  Which form is shorter
depends on the case (for the cyclic quintic the Fourier form has 140 leaves in Wolfram's
count, the eigenvector form 173; for cubics both forms coincide with Cardano's).  The
implementations compute both and keep the shorter, or use the one requested.

\subsection{Branch selection}\label{sec:branch}

The candidates $\zeta^e Q^{1/q}$, $e=0,\dots,q-1$, for $Q=R_k^q\neq0$ are the $q$ distinct
$q$-th roots of $Q$; they are separated by at least
$|Q|^{1/q}\cdot2\sin(\pi/q)$.  The value of $R_k$ is known independently from the exact
coordinates of $R_k$ in the splitting field (Section~\ref{sec:engine}), so a numerical
comparison at sufficient precision identifies $e$.  In Wolfram the comparison is made at
60 digits with a separation check, and the final identity is verified exactly by
\code{RootReduce}; in Python it is made in ball arithmetic and is rigorous by itself: a
candidate is accepted only if it is the unique one whose ball overlaps the ball of $R_k$
(the true candidate always overlaps, so uniqueness identifies it), and the precision is
doubled until this happens.

\emph{The branch cut.}  There is one trap in rigorous evaluation that none of the nine
reports mentions.  When $Q$ is a real negative number that is \emph{written} as a sum of
complex-conjugate radicals (which happens as soon as a real number is expressed through
non-real radicands, for instance every real root of a cubic with three real roots), its
computed enclosure is a small box straddling the negative real axis.  The principal
$q$-th root is discontinuous there, so the enclosure of $Q^{1/q}$ covers both limits and
no precision separates the candidates.  The exact data resolve this: $Q$ is real if and
only if it is fixed by the automorphism that acts as complex conjugation on the roots (a
permutation the engine identifies), and its sign is then certified from the real part.
In that case the implementation writes the principal root as $(-1)^{1/q}(-Q)^{1/q}$,
an identity of principal branches for real negative $Q$, and the enclosure of $-Q$ lies
near the positive axis where the root is continuous.  The same rewriting is applied in
the structural layer to discriminants and to binomial right-hand sides, whose exact
values are known as algebraic numbers.

\subsection{The algorithm and its completeness}

\begin{algorithm}\label{alg:descent}
Input: $a=\Root[p,k]$ with $p$ irreducible.
\begin{enumerate}[nosep]
\item Compute $G$ as a permutation group on the roots, with the splitting field $L$ in
coordinates (Section~\ref{sec:engine}).  If $G$ is not solvable, stop: no radical
expression exists.
\item Let $m$ be the product of the odd primes dividing $|G|$.  Compute the same data for
$p\cdot\prod_{q\mid m}\Phi_q$ (omitting a factor equal to $p$); this is $M=L(\zeta_m)$
with its group.  Locate $\zeta_q$ among the roots and let $H$ be the stabilizer of all of
them.
\item Build a composition series of $H$ with prime quotients, a generator $\sigma_i$ of
each quotient, and the coordinates of $a$.
\item Define $\rho(v,i)$ for $v\in F_i$ recursively: if $i=0$, solve for the coordinates of
$v$ in the basis $\{\prod_q\zeta_q^{e_q}:0\le e_q\le q-2\}$ of $B$ and return the
corresponding expression in $(-1)^{2/q}$; if $v$ is fixed by $H_{i-1}$, return
$\rho(v,i-1)$; otherwise form the $R_k$ of Proposition~\ref{prop:lagrange}, compute
$R_k^q$ and the $c_k$ by exact linear algebra in coordinates, select branches as in
Section~\ref{sec:branch}, and return \eqref{eq:fourier} or \eqref{eq:eigen} with
$\rho(\cdot,i-1)$ applied to $R_0$, $R_k^q$ and $c_k$.  Memoize on $(v,i)$.
\item Output $\rho(a,r)$ and verify $\rho(a,r)-a=0$ exactly.
\end{enumerate}
\end{algorithm}

\begin{theorem}[Completeness]\label{thm:complete}
Algorithm~\ref{alg:descent} terminates.  If $G$ is solvable it outputs a radical
expression equal to $\Root[p,k]$ of depth at most $r+1$, where $r$ is the length of the
composition series of $H$; otherwise it reports that no radical expression exists.
\end{theorem}

\begin{proof}
Nonsolvability is decided exactly by the commutator series.  If $G$ is solvable, $H$ is
solvable by Lemma~\ref{lem:base}(ii) and the series exists; by Lemma~\ref{lem:base}(iii),
$\zeta_{q_i}\in B\subseteq F_{i-1}$ for every $i$, so Proposition~\ref{prop:lagrange}
applies at every level.  The recursion $\rho(v,i)$ calls $\rho$ only with second argument
$i-1$ and terminates.  Correctness: at level $0$ the element lies in $B$ and its
expression in the power basis of the $\zeta_q$ is exact; at level $i$, the identity
\eqref{eq:fourier} or \eqref{eq:eigen} holds exactly in $M$ with the true branches, and
the true branches are the selected ones because the candidates are separated and the
comparison is made at a precision exceeding the separation (rigorously so in ball
arithmetic).  Each level adds at most one nesting of radicals to the expressions coming
from the level below, and the base expressions have depth one, so the depth is at most
$r+1$.  Finally, the exact verification in step 5 is an independent check.
\end{proof}

\begin{remark}[Cost]
The cost is dominated by step 2: the splitting field of $p\prod\Phi_q$ has degree
$|H|\cdot\varphi(m)$ over $\Q$, and our engine builds it by adjoining one root at a
time with \code{RootReduce}; for the two examples of the question the extended groups
have orders $48$ and $40$, computed in about $33$ and $10$ seconds in Wolfram.  The
group order limit (\code{"MaxGroupOrder"}, default $400$) turns larger cases into an
explicit \code{ResourceLimit} status rather than a wrong answer.  The only heuristic
component is the choice of random weights in the resolvent construction, which is
checked by exact counting and closure tests.
\end{remark}

\subsection{Two refinements}\label{sec:refinements}

\emph{The tensor algebra.}  Several reports adjoin $\zeta_q$ by working in
$L[z]/\Phi_q(z)\cong L\otimes_\Q\Q(\zeta_q)$ instead of in the field $L(\zeta_q)$.  This
is a ring, not a field, as soon as $\zeta_q\in L$ ($\Phi_q$ splits in $L$), and it has
zero divisors; report~05 observes this and makes its descent division-free.  Our
implementations avoid the issue by computing the splitting field of $p\prod\Phi_q$
directly, which is a field whose degree is the true $[M:\Q]$ (for $x^3-2$, adjoining
$\Phi_3$ does not enlarge the field, and the engine sees a group of order $6$, not $12$).
The one division we perform, $c_k=R_k/R_{k_1}^m$, is in a field and is exact.

\emph{Fourier versus eigenvector.}  The Fourier form \eqref{eq:fourier} needs no
nonzero-generator search (report~03), the eigenvector form \eqref{eq:eigen} usually
gives shorter formulas (report~04).  The two forms are equally valid and we implement
both; the eigenvector form needs the extra exactness argument of
Proposition~\ref{prop:lagrange}(iii), which we have supplied.

\section{Negative results without the group}\label{sec:negative}

Computing $G$ exactly is the expensive step (for $x^5-x-1$, whose group is $S_5$ of
order $120$, our engine needs $18$ minutes in Wolfram).  Two classical facts give
nonsolvability proofs in milliseconds.

\begin{theorem}[Dedekind--Frobenius]\label{thm:frobenius}
Let $\ell$ be a prime not dividing the discriminant or the leading coefficient of $p$.
If $p\bmod\ell$ factors into irreducible factors of degrees $d_1,\dots,d_s$, then $G$
contains a permutation of cycle type $(d_1,\dots,d_s)$.
\end{theorem}

\begin{theorem}[Galois]\label{thm:agl}
A solvable transitive permutation group of prime degree $n$ is conjugate to a subgroup
of $\AGL(1,n)=\{x\mapsto ax+b\}$ acting on $\F_n$.  Consequently each of its elements has
cycle type $1^n$, $n$, or $1\,d^{(n-1)/d}$ for some divisor $d$ of $n-1$.
\end{theorem}

\begin{proof}
The first statement is Galois' theorem on solvable transitive groups of prime degree
\cite[Chapter 4]{Milne}, \cite[\S 14.7]{DummitFoote}.  For the second, $x\mapsto x+b$ with
$b\ne0$ is an $n$-cycle; $x\mapsto ax+b$ with $a\ne1$ has the unique fixed point
$x=b/(1-a)$ and, after conjugating by a translation, is $x\mapsto ax$, whose nonzero
orbits all have length $\operatorname{ord}(a)\mid n-1$.
\end{proof}

\begin{corollary}\label{cor:frob}
If $n$ is prime and some Frobenius cycle type is not of the form in
Theorem~\ref{thm:agl}, then $G$ is not solvable and $\Root[p,k]$ has no radical
expression.
\end{corollary}

For example $x^5-x-1\bmod 3$ has factors of degrees $(2,3)$ and $x^7-x-1\bmod 3$ has
degrees $(1,2,4)$; neither is admissible, so both are settled immediately.

For composite degrees we use a second criterion.

\begin{proposition}\label{prop:longcycle}
Let $n$ be composite and not a prime power.  If $G$ contains a cycle of prime length
$\ell>n/2$ (with all other points fixed), then $G$ is not solvable.
\end{proposition}

\begin{proof}
A transitive group containing an $\ell$-cycle with $\ell>n/2$ is primitive: if $\Delta$
were a block system with blocks of size $1<b<n$, the $\ell$-cycle would permute the
blocks it meets in a cycle of length dividing $\ell$; length $1$ would force its support
of size $\ell$ into a single block of size $b\le n/2<\ell$, and length $\ell$ would need
$\ell b\le n$, impossible for $b\ge2$.  A solvable primitive group has a unique minimal
normal subgroup, which is elementary abelian and regular, so its degree is a prime power
\cite[Chapter II]{Wielandt}.  Since $n$ is not a prime power, $G$ is not solvable.
\end{proof}

Both criteria are inconclusive in the other direction (they never prove solvability),
and neither applies to composite prime-power degrees $4,8,9,16,\dots$; in those cases the
implementations compute the group exactly, subject to the group order limit.  For
$x^6-x-1$ (degree $6$, group $S_6$) and $x^6+x^4+3x^2-2x+5$ (group of order $120$) a
$5$-cycle is found within the first few primes and the Frobenius test proves
nonsolvability in about $10$ milliseconds (a full \code{RootToRadicals} call takes up to
a second more, spent in the built-in \ToRadicals{} attempt that precedes it).

\section{What the notebook of the question does}\label{sec:notebook}

The notebook \code{FindExtension.nb} attached to the question (an inert transcription
of its input cells is kept in \code{reports/report-09/original/}) is not about the two
examples above.  Its target is
\[
\beta_0=I^{-1}_{1/9}\bigl(\tfrac13,\tfrac13\bigr),
\]
the inverse of the regularized incomplete beta function, which the author evaluates to
$1000$ digits and recognizes with \code{RootApproximant} as an algebraic number; the
notebook then tries to express that number through nested radicals and $\cos(\pi/9)$.
We reproduced the recognition: \code{RootApproximant} returns in $7$ seconds a monic
integer polynomial of degree $27$ (leading terms $x^{27}+108x^{26}-2970x^{25}+\dots$,
constant term $1$), the identity holds to $400$ digits, and over $\Q(3^{1/6})$ the
polynomial factors into pieces of degrees $9$ and $18$, which is the tower the notebook
finds.  Since $27=3^3$ is a prime power, neither Frobenius test of
Section~\ref{sec:negative} applies, and the exact group of the degree-$27$ polynomial
did not finish within $15$ minutes in Wolfram even with the order limit raised to
$3000$.  The notebook's tower settles the question nevertheless: the degree-$9$ number
$\theta$ that the notebook adjoins (the third root of
$x^9-657x^8+6111x^7+3318x^6+19647x^5-12033x^4+3972x^3-684x^2+9x-1$) has a Galois group
of order $18$, and \code{RootToRadicals} expresses it in radicals by two Kummer steps
of prime $3$ in $43$ seconds (verified exactly; depth $3$, $111$ leaves); the degree-$9$
factor of the minimal polynomial over $\Q(3^{1/6})$ that vanishes at $\beta_0$ splits
over $\Q(3^{1/6},\theta)$ into three cubics ($237$ seconds), so $\beta_0$ is a root of
a cubic with radical coefficients and therefore has a radical expression.  We have not
assembled that expression.  The notebook's devices, in the order they appear, are:

\begin{description}[nosep]
\item[\code{radicalDepth}] the nesting depth of rational powers (and of \code{Sin},
\code{Cos}), the same measure as our \code{RadicalDepth}.
\item[\code{findExtension}] eliminates $p,q$ from
$\code{PolynomialRemainder}[mp,\,x^2+px+q]=0$ and returns the simplest $p$ with a
solution.  A monic quadratic dividing the minimal polynomial has $-p=a_i+a_j$, so
this is a search over the pair sums, i.e.\ over the roots of the pair-sum resolvent
$R_2$ of Lemma~\ref{lem:pairsum}, with \code{Simplify`SimplifyCount} as the
selection criterion.  \code{findExtension3} does the same with cubic factors
$x^3+px^2+qx+r$, i.e.\ with triple sums.  The systematic replacement is
Proposition~\ref{prop:pairsum}: compute $R_2$ exactly, factor it over $\Q$, and factor
$p$ over $\Q(y_0)$ for a low-degree root $y_0$.
\item[\code{project}] tests whether the real or imaginary part of some conjugate has
smaller degree than the number itself.  This is a special case of the same idea:
$2\operatorname{Re}(a)=a+\bar a$ is a pair sum.
\item[\code{factor}] factors the minimal polynomial over an extension and picks the
factor vanishing at the number with \code{PossibleZeroQ}; this is the factorization
step of Proposition~\ref{prop:pairsum}.
\item[\code{solve}] picks, among the \code{Solve} roots of a low-degree factor, the
one equal to the number (\code{PossibleZeroQ} again): branch selection
(Section~\ref{sec:branch}) done numerically.
\item[\code{fuzz}] draws random subsets of the conjugates, sums them, guesses a
low-degree algebraic number with \code{RootApproximant} and checks it with
\code{PossibleZeroQ}.  This is a randomized search for elements of proper subfields of
the splitting field; the Galois-theoretic replacement is the list of subgroups and
fixed fields that the engine computes, or, for the specific purpose of radicals, the
composition series of Section~\ref{sec:descent}, which needs no search at all.
\item[\code{simplify}] iterates \code{FullSimplify} with a custom complexity function
that penalizes \code{ArcTan}, hyperbolic functions and every \Root{} object except
one chosen degree-$9$ number.  This is the only device without a systematic
replacement: the shortness of a radical expression is not a Galois-theoretic
invariant, and the reports that discuss it (04, 09) offer heuristics only.
\end{description}

The notebook's approach is sound as a discovery procedure and it does find the tower
$\Q\subset\Q(3^{1/6})\subset\Q(3^{1/6},\theta)$ over which the target becomes a root of
a cubic; what it lacks is a termination guarantee and a criterion for when to stop
looking.  The descent of Section~\ref{sec:descent} supplies both, at the price of the
splitting field, and for the piece $\theta$ that the notebook could not simplify it
produces the radical expression directly.

\section{Comparison of the nine reports}\label{sec:reports}

Nine reports on this question were written independently before this article; they are
kept verbatim in \code{reports/report-01} to \code{report-09}.  All nine share the
same core: the criterion on the minimal polynomial, the cyclotomic base with
$m=\rad|G|$, a prime composition series, Lagrange or Fourier resolvents, exact branch
selection, the same two closed forms for the examples, and a completeness theorem
conditional on exact primitives.  Table~\ref{tab:reports} lists what each one adds.

\begin{table}[htbp]
\centering\small
\begin{tabularx}{\textwidth}{@{}c>{\raggedright\arraybackslash}X>{\raggedright\arraybackslash}X@{}}
\toprule
report & distinctive contribution & what was executed\\
\midrule
01 & Kummer generators as a rational nullspace; ``rational rectangle'' embeddings
with exact sign certificates. & The only forced general run on the two original
examples; neither finished within its 480-second limit (the sextic reached the degree-48
field, the quintic the degree-40 field).  Two prime-5 Kummer steps (cyclic and binomial
quintic, 22 and 28 s).\\
02 & The pair-sum resolvent $\Res_x(f(x),f(y-x))=2^nf(y/2)R_2(y)^2$ identified as the
systematic form of the notebook's \code{findExtension}; a trace--Fourier basis search
with a termination proof; a native Wolfram \code{kummerSolve}. & Structural paths on the
examples (0.7--2.6 s per root); forced descent only on $x^3-2$, $x^4-2$.\\
03 & Fourier inversion $a=\frac1p\sum_jR_j$ needing no generator search; the clearest
account of SymPy embedding pitfalls; the point that an oversized supplied field proves
nothing about nonsolvability; a native Wolfram \code{galoisConvert}. & Structural paths
on the examples (1.8 s, 1.2 s).\\
04 & Relative-automorphism recovery without a second factorization; trace-projector
scoring for small radicands; an independent arithmetic-tower checker. & Structural paths
on the examples (0.7 s, 0.2 s); the cyclic quintic by descent (16.7 s).\\
05 & Division-free tensor algebra $L\otimes\Q(\zeta_p)$ with zero divisors (with a test
multiplying two of them); resultant primitive elements with a square-free certificate;
norm-polynomial justification of branch comparisons. & Structural paths (0.6 s, 0.2 s);
the cyclic quintic by Fourier descent (1.0 s); a forced $S_4$ run hit its resource limit
at 57 s.\\
06 & Exact Frobenius negative test; a principal-sector branch predicate with
conjugation-based zero tests; straight-line register programs as certificates. &
Structural paths on all conjugates of both examples (up to 104 s for one sextic root);
the forced cyclic quintic timed out at 90 s.\\
07 & A native Wolfram \code{SystematicRadicals.wl} with Kummer eigenprojections; a
notebook audit with line numbers. & Structural paths (2.3 s, 1.6 s); its forced run on
$x^5-2$ was inconclusive at 30 s.\\
08 & The coherent-branch theorem (any consistent branch assignment reproduces the whole
root set); a purely rational certificate checker that proves the root set rather than an
index. & Structural paths (milliseconds for the selected roots).\\
09 & A five-status taxonomy (\code{Success}, \code{NotSolvable}, \code{SearchExhausted},
\code{ResourceLimit}, \code{VerificationFailure}); replayable certificates with an
embedding anchor; a line-numbered notebook critique. & Structural paths (4.0 s, 0.3 s);
forced descents on $x^5-2$ (34 s) and a nonabelian relative group (39 s).\\
\bottomrule
\end{tabularx}
\caption{What each of the nine reports adds to the common core, and what its Python code
actually executed (its own reported timings).}\label{tab:reports}
\end{table}

All nine reports state the same two closed forms for the examples (report~08 writes the
sextic with two positive cube roots, $((9+\sqrt{849})/18)^{1/3}-((\sqrt{849}-9)/18)^{1/3}$,
which is the same number).  Three reports (02, 03, 07) ship a native Wolfram
implementation of the general algorithm; none of the three was executed by its author.
The only report that forced the general descent on the original examples (01) did not
finish, because it built the splitting field in the power basis of a primitive element
(Section~\ref{sec:engine}).  Four reports (01, 04, 05, 09) executed a prime-5 Kummer step
on some input.  The status taxonomy of report~09 is adopted here, with
\code{SearchExhausted} renamed \code{NotFound}.

Three points needed correction or independent verification.  Report~07 cites Milne
``Theorem 3.28'' for the solvability criterion; the theorem is in Chapter 5 of
\cite{Milne}.  Several reports refer to \code{MANIFEST} or \code{SHA256SUMS} files that
are not present here (checksum ledgers are not kept in this repository).  And no report
executed any Wolfram Language code: every Wolfram claim in them was static.  We ran the
Wolfram side natively; the results are in Section~\ref{sec:measurements}.

Every report proposes to build the splitting field with
\code{ToNumberField[roots, All]} (or its SymPy analogue \code{primitive\_element}).
The companion project measured that call at $14$ minutes for a degree-$36$ field, and
the reason is structural: the power basis of a generic primitive element has
coefficients of enormous height.  Our implementations use the tower basis with
trace-form coordinates of the companion engine instead (Section~\ref{sec:engine}), which
is why the forced descent on the examples takes seconds rather than failing to finish, as
it did in report~01.

\section{Algorithms and implementation}\label{sec:impl}

\subsection{The splitting-field engine}\label{sec:engine}

Both implementations reuse the Galois engine of the companion project
(\code{RootDecomposition.wl} and \code{rootdecomp.py} in \path{root-decomposition/}).  It
computes the Galois group by numerical resolvents: roots are adjoined one at a time,
$\theta_{k}=\theta_{k-1}+w_ka_k$ with small random integer weights, the exact minimal
polynomial of $\theta_k$ is found with \code{RootReduce} (Wolfram) or by rigorous
rounding of a product of balls (Python), and the orbit of partial tuples is matched to
its roots numerically with counting checks; the result is verified by group closure.  The
splitting field is represented in the tower monomial basis $\prod_j a_{g_j}^{e_j}$, and
the coordinates of an element are obtained from its numerical conjugates through the
trace form (the Gram matrix of the basis is an integer matrix), so multiplication
matrices, automorphism matrices and fixed fields are exact rational data.  Elements of
$M$ are therefore rational vectors, $R_k$, $R_k^q$ and $c_k$ are computed by exact
matrix arithmetic, and the numerical value of any element is available at any precision.

\subsection{Interface}

\begin{lstlisting}
Get["RootToRadicals.wl"];
RootToRadicals[Root[-1 - #^2 - #^3 + #^4 + #^6 &, 2]]
RootRadicalReport[Root[6 + 25 # - 25 #^3 + 5 #^5 &, 5]]
RootRadicalReport[Root[#^5 + #^4 - 4 #^3 - 3 #^2 + 3 # + 1 &, 1], Method -> "Galois"]
RootSolvableQ[Root[#^5 - # - 1 &, 1]]         (* False, by Frobenius cycle types *)
\end{lstlisting}
\begin{sloppypar}
\code{RootToRadicals[a]} returns a radical expression or a \code{Failure} whose tag is one
of \code{NotSolvable} (proved), \code{NotFound} (structural search exhausted, only with
\code{Method -> "Structural"}), \code{ResourceLimit}, \code{VerificationFailed},
\code{Inexact}, \code{NotAlgebraic}, \code{InvalidOptions}.  \code{RootRadicalReport}
returns an association with the expression, the method (\code{ToRadicals},
\code{Decompose}, \code{Reciprocal}, \code{Dickson}, \code{PairSum}, \code{Galois}), the
verification status (\code{True} by \code{RootReduce}, or \code{Indeterminate} when the
exact check exceeded its time limit and only a 200-digit numerical check passed), the
depth, the leaf count, the group orders and the primes of the composition series, and
the timing.  Options: \code{Method} (\code{Automatic}, \code{"Structural"},
\code{"Galois"}), \code{"Resolvents"} (\code{Automatic}, \code{"Fourier"},
\code{"Eigenvector"}), \code{"MaxGroupOrder"}, \code{"WorkingPrecision"},
\code{"VerificationTimeLimit"}, \code{"MaxDepth"}.  \code{RadicalExpressionQ} and
\code{RadicalDepth} implement the grammar of Section~\ref{sec:conventions}.
\end{sloppypar}

\begin{lstlisting}
import roottoradicals as rt
r = rt.root_to_radicals("Root[-1 - #^2 - #^3 + #^4 + #^6 &, 2]")
r.expression, r.method, r.verified, r.wolfram()
rt.root_to_radicals("Root[1 + 3 # - 3 #^2 - 4 #^3 + #^4 + #^5 &, 1]", method="galois")
rt.is_solvable("Root[-1 - # + #^5 &, 1]")      # False
\end{lstlisting}
The Python function returns a \code{RadicalResult} (a SymPy expression with the same
metadata) or raises \code{NotSolvable}, \code{NotFound}, \code{ResourceLimit} or
\code{PrecisionError}.  Its structural layer has the classical formulas for degree
$\le4$, \code{Decompose}, reciprocal symmetry and Dickson polynomials; the pair-sum
reduction, which needs factorization over a number field, is Wolfram-only, and the
general descent covers those cases in Python.  \code{verified} is a rigorous ball check
that the expression encloses the requested root and no other conjugate;
\code{python/verify\_wolfram.py} additionally verifies Python's expressions exactly in a
Wolfram kernel with \code{RootReduce}.

\subsection{Measurements}\label{sec:measurements}

All timings were measured on one Windows machine, in Wolfram 15.0.1 (\code{wolfram -script})
and in Python 3.14 with python-flint 0.8 and SymPy 1.14, with the default options
(\code{"MaxGroupOrder"} $400$, 80 digits or 300 bits).  ``Orders'' are $|G|$ and the
order of the group of $p\prod\Phi_q$.

\begin{center}
\small
\begin{tabular}{@{}llrrr@{}}
\toprule
input & method & orders & Wolfram & Python\\
\midrule
sextic example $\alpha$ & Reciprocal & & 0.12 s & 0.07 s\\
quintic example $\beta$ & Dickson & & 0.04 s & 0.01 s\\
sextic example, descent forced & Galois & $24\to48$ & 41 s & 24 s\\
quintic example, descent forced & Galois & $20\to40$ & 18 s & 12 s\\
cyclic quintic $x^5+x^4-4x^3-3x^2+3x+1$ & Galois & $5\to20$ & 8.4 s & 0.37 s\\
\quad the same, Fourier form only & Galois & $5\to20$ & 4.8 s & 0.05 s\\
$x^4-x-1$, descent forced & Galois & $24\to48$ & 32 s & 12 s\\
$x^5-2$, descent forced & Galois & $20\to20$ & 4.8 s & 1.6 s\\
$\Phi_7$, descent forced & Galois & $6\to12$ & 1.5 s & 0.24 s\\
cyclic quintic composed with $x^2$ (degree 10) & Decompose & $5\to20$ & 9.3 s & 0.19 s\\
$D_7(x,1)-3$ & Dickson & & 0.10 s & 0.02 s\\
sextic with pair sums of degree 3 & PairSum / Galois$^*$ & $12\to24$ & 0.19 s & 0.83 s\\
$x^5-x-1$: not solvable & Frobenius & & 1.2 s & 0.26 s\\
$x^6-x-1$: not solvable & Frobenius & & 0.01 s & 0.01 s\\
$x^5-x-1$ by the exact group (order 120) & group & $120$ & 1088 s & ---\\
\bottomrule
\end{tabular}
\end{center}
\noindent{\small $^*$ The pair-sum reduction is Wolfram-only; Python solves this input by the descent.}

The Wolfram times for the descent are dominated by the construction of the extended
splitting field (about $33$ s of the $41$ s for the sextic); Python's Arb-based engine is
faster there.  The Wolfram time for the cyclic quintic with the default
\code{"Resolvents" -> Automatic} includes computing both forms; the eigenvector form is
longer for this input (173 versus 140 leaves) and is discarded.  Wolfram's structural
layer uses the built-in \code{ToRadicals} first (which handles degree $\le4$, binomials
and decomposable inputs in milliseconds), so its \code{Decompose} path is only reached
for compositions whose outer piece the built-in function cannot solve, as in the
degree-10 example.

The regression suites were executed: \code{RunTests.wl} runs $60$ tests in $90$ s of
wall time (including the forced descents on both examples and all their conjugates);
\code{test\_roottoradicals.py} runs $46$ cases in $51$ s.  \code{verify\_wolfram.py}
re-verifies $13$ Python expressions exactly in a Wolfram kernel with \code{RootReduce};
all pass.

\subsection{Limits and scope of the certificates}

\begin{itemize}[nosep]
\item A returned expression is exactly equal to the input: in Wolfram by
\code{RootReduce}; in Python by the exactness of the coordinate identities together with
rigorous branch selection, and again by the final ball check.  The status
\code{Verified -> Indeterminate} (Wolfram) is reported honestly when the exact check hits
its time limit.
\item \code{NotSolvable} is a proof: either a Frobenius cycle type
(Corollary~\ref{cor:frob}, Proposition~\ref{prop:longcycle}) or the exact commutator
series of the computed group.  The group computation relies on numerical matching with
counting and closure checks; its exactness has the same standing as in the companion
project.
\item \code{ResourceLimit} asserts nothing about solvability.  Nonsolvable polynomials of
composite prime-power degree with a group larger than the limit end here; so do solvable
ones with $|H|\varphi(m)$ beyond the limit.
\item The structural layer never asserts nonexistence; \code{NotFound} is returned only
when the general descent was disabled.
\item Root indices of non-real roots are exchanged between Wolfram and Python by
value, not by index: Python orders conjugate pairs by real part and imaginary
magnitude, and Wolfram's ordering of non-real roots can depend on its isolation method.
\end{itemize}

\section{Conclusion}

The restriction of \ToRadicals{} is one of scope, not of principle.  A radical expression
exists exactly when the Galois group of the minimal polynomial is solvable, which is
decidable in polynomial time; constructing the expression may require the splitting
field, and this is where a general library function stops.  With the splitting field
available in exact coordinates, the Galois--Kummer descent is a short, fully proved,
complete algorithm, and the structural reductions give the short formulas one wants in
practice, including the two of the question.  Both are implemented, tested and measured
in two systems.

\begin{thebibliography}{9}
\bibitem{LandauMiller} S.~Landau and G.~L.~Miller, \emph{Solvability by radicals is in
polynomial time}, J. Comput. System Sci. 30 (1985), 179--208.
\bibitem{Milne} J.~S.~Milne, \emph{Fields and Galois Theory}, version 5.10, 2022,
\url{https://www.jmilne.org/math/CourseNotes/ft.html}.
\bibitem{DummitFoote} D.~S.~Dummit and R.~M.~Foote, \emph{Abstract Algebra}, 3rd ed.,
Wiley, 2004.
\bibitem{Wielandt} H.~Wielandt, \emph{Finite Permutation Groups}, Academic Press, 1964.
\bibitem{LidlMullenTurnwald} R.~Lidl, G.~L.~Mullen and G.~Turnwald, \emph{Dickson
Polynomials}, Pitman Monographs 65, Longman, 1993.
\bibitem{Cohen} H.~Cohen, \emph{A Course in Computational Algebraic Number Theory},
Springer, 1993.
\bibitem{Question} Mathematica StackExchange question 34011,
\emph{Why is ToRadicals not able to handle all cases? Is there a workaround?},
\url{https://mathematica.stackexchange.com/q/34011}.
\end{thebibliography}

\end{document}
